% Generated by tools/paper/render.py from main_text.py; edit the source there, not this file.
\section{Discussion}

Phasers are usually compared through their phased variants, but many uses of long-read phasing depend on the reads: a read left unphased is lost to every haplotype-resolved analysis, and a block joined in the wrong orientation misassigns all of its reads. On these terms pgphase phased more reads correctly than the tools we compared it with, on every chromosome we tested. Two of its design choices are the likely reasons, although we have not measured their contributions separately. First, it has more sites to phase reads with: the graph's known alleles, observed in the reads' graph alignments rather than called, and the sample's own heterozygous variants. Second, it is conservative in joining blocks: sites at collapsed duplications, poorly mapped sites and homopolymer length calls still label the reads that cross them but cannot join two blocks. The cost of the second is more and shorter blocks than HiPhase produces where the reads cannot verify how blocks connect.

Our evaluation has limits. All tools were given reads aligned to a pangenome graph and surjected onto CHM13, not the linear alignments (for example from minimap2 or pbmm2) the other tools are usually run on, and their accuracy with such alignments may differ. We evaluated one sample and one sequencing technology on eight chromosomes; other samples, ONT reads and other graphs remain to be tested. The development chromosomes shaped pgphase's rules and thresholds, and although the chromosomes not used for tuning showed a similar margin, they were scored during development. The read truth relies on alignment to the parental assemblies, and its labels are least reliable where the parental haplotypes are most similar.

pgphase also has limits of its own. Its catalog holds only the alleles of the genomes in the graph: at a tandem repeat whose expansion in the sample no genome shares, a read's graph alignment follows the closest known walk, and most of the switches we could not classify are of this kind. Representing such alleles from the reads themselves is a natural next step. pgphase also needs graph alignments of the reads and more memory than the phasers it was compared with. Finally, a record-level comparison counts differences of representation in repeats as switch errors; classifying switches by haplotype sequence, as here, or comparing haplotypes rather than records, would make phasers easier to compare.
